\begin{env}[4.4.3]
\label{0.4.4.3-fr}
Vamos a ver ahora que se puede definir un \emph{isomorfismo}
canónico de \emph{bifuntores} en $\mathcal{F}$ y $\mathcal{G}$
\[
  \operatorname{Hom}_{\mathcal{A}}(\Psi^*(\mathcal{G}),\mathcal{F})
  \xrightarrow{\sim}
  \operatorname{Hom}_{\mathcal{B}}(\mathcal{G},\Psi_*(\mathcal{F}))
  \tag{4.4.3.1}\label{0.4.4.3.1-fr}
\]
que designaremos por $v\to v_\theta^\flat$ (o simplemente
$v\to v^\flat$ si no es posible ninguna confusión); denotaremos por
$u\to u_\theta^\sharp$, o $u\to u^\sharp$ el isomorfismo inverso.
Esta definición es la siguiente: al componer
$v:\Psi^*(\mathcal{G})\to\mathcal{F}$ con la aplicación canónica
$\psi^*(\mathcal{G})\to\Psi^*(\mathcal{G})$, se obtiene un homomorfismo de
haces de grupos $v':\psi^*(\mathcal{G})\to\mathcal{F}$, que es también un
homomorfismo de $\psi^*(\mathcal{B})$-módulos. Se deduce de él (3.7.1) un
homomorfismo
${v'}^\flat:\mathcal{G}\to\psi_*(\mathcal{F})=\Psi_*(\mathcal{F})$, que es
también un homomorfismo de $\mathcal{B}$-Módulos como se
\oldpage[0\textsubscript{I}]{43}
verifica sin dificultad; se toma $v_\theta^\flat={v'}^\flat$. Del mismo modo, de
$u:\mathcal{G}\to\Psi_*(\mathcal{F})$, que es un homomorfismo de
$\mathcal{B}$-Módulos, se deduce (3.7.1) un homomorfismo
$u^\sharp:\psi^*(\mathcal{G})\to\mathcal{F}$ de
$\psi^*(\mathcal{B})$-Módulos, de donde, por tensorización con $\mathcal{A}$, un
homomorfismo de $\mathcal{A}$-Módulos
$\Psi^*(\mathcal{G})\to\mathcal{F}$, que designaremos por
$u_\theta^\sharp$. Es inmediato verificar que
$(u_\theta^\sharp)_\theta^\flat=u$ y
$(v_\theta^\flat)_\theta^\sharp=v$, así como el carácter funtorial en
$\mathcal{F}$ del isomorfismo $v\to v_\theta^\flat$. El carácter
funtorial en $\mathcal{G}$ de $u\to u_\theta^\sharp$ se deduce entonces
formalmente como en (3.5.4) (razonamiento que demostraría también el
carácter funtorial de $\Psi^*$ establecido directamente en (4.3.1)).

Si se toma como $v$ el homomorfismo identidad de $\Psi^*(\mathcal{G})$,
$v_\theta^\flat$ es un homomorfismo
\[
  \rho_{\mathcal{G}}:\mathcal{G}\longrightarrow
  \Psi_*(\Psi^*(\mathcal{G}))~;
  \tag{4.4.3.2}\label{0.4.4.3.2-fr}
\]
si se toma como $u$ el homomorfismo identidad de $\Psi_*(\mathcal{F})$,
$u_\theta^\sharp$ es un homomorfismo
\[
  \sigma_{\mathcal{F}}:\Psi^*(\Psi_*(\mathcal{F}))
  \longrightarrow\mathcal{F}~;
  \tag{4.4.3.3}\label{0.4.4.3.3-fr}
\]
estos homomorfismos se llamarán \emph{canónicos}. En general no son ni
inyectivos ni sobreyectivos. Se tienen factorizaciones canónicas análogas a
(3.5.3.3) y (3.5.4.4).

Se observará que si $s$ es una sección de $\mathcal{G}$ sobre un abierto
$V$ de $Y$, $\rho_{\mathcal{G}}(s)$ es la sección $s'\otimes1$ de
$\Psi^*(\mathcal{G})$ sobre $\psi^{-1}(V)$, siendo $s'$ tal que
$s'_x=s_{\psi(x)}$ para todo $x\in\psi^{-1}(V)$. Observemos también que si
$u:\mathcal{G}\to\psi_*(\mathcal{F})$ es un homomorfismo, define para
todo $x\in X$ un homomorfismo
$u_x:\mathcal{G}_{\psi(x)}\to\mathcal{F}_x$ sobre los tallos, obtenido al
componer
$(u^\sharp)_x:(\Psi^*(\mathcal{G}))_x\to\mathcal{F}_x$ y el homomorfismo
canónico $s_x\to s_x\otimes1$ de $\mathcal{G}_{\psi(x)}$ en
$(\Psi^*(\mathcal{G}))_x
=\mathcal{G}_{\psi(x)}\otimes_{\mathcal{B}_{\psi(x)}}\mathcal{A}_x$.
El homomorfismo $u_x$ se obtiene también por paso al límite inductivo a
partir de los homomorfismos
\[
  \Gamma(V,\mathcal{G})\xrightarrow{u}
  \Gamma(\psi^{-1}(V),\mathcal{F})\longrightarrow\mathcal{F}_x,
\]
donde $V$ recorre las vecindades de $\psi(x)$.
\end{env}

\begin{env}[4.4.4]
\label{0.4.4.4-fr}
Sean $\mathcal{F}_1$, $\mathcal{F}_2$ unos $\mathcal{A}$-Módulos,
$\mathcal{G}_1$, $\mathcal{G}_2$ unos $\mathcal{B}$-Módulos, $u_i$ ($i=1,2$)
un homomorfismo de $\mathcal{G}_i$ en $\mathcal{F}_i$. Designaremos por
$u_1\otimes u_2$ el homomorfismo
\[
  u:\mathcal{G}_1\otimes_{\mathcal{B}}\mathcal{G}_2
  \longrightarrow
  \mathcal{F}_1\otimes_{\mathcal{A}}\mathcal{F}_2
\]
tal que $u^\sharp=(u_1)^\sharp\otimes(u_2)^\sharp$ (teniendo en cuenta
(4.3.3.1)); se verifica que $u$ es también el compuesto
\[
  \mathcal{G}_1\otimes_{\mathcal{B}}\mathcal{G}_2
  \longrightarrow
  \Psi_*(\mathcal{F}_1)\otimes_{\mathcal{B}}\Psi_*(\mathcal{F}_2)
  \longrightarrow
  \Psi_*(\mathcal{F}_1\otimes_{\mathcal{A}}\mathcal{F}_2),
\]
donde la primera flecha es el producto tensorial ordinario
$u_1\otimes_{\mathcal{B}}u_2$ y la segunda el homomorfismo canónico
(4.2.2.1).
\end{env}

\begin{env}[4.4.5]
\label{0.4.4.5-fr}
Sean $(\mathcal{G}_\lambda)_{\lambda\in L}$ un sistema inductivo de
$\mathcal{B}$-Módulos, y, para todo $\lambda\in L$, sea $u_\lambda$ un
homomorfismo $\mathcal{G}_\lambda\to\Psi_*(\mathcal{F})$, que forman un sistema
inductivo; pongamos
$\mathcal{G}=\varinjlim\mathcal{G}_\lambda$ y
$u=\varinjlim u_\lambda$; entonces los $(u_\lambda)^\sharp$ forman un sistema
inductivo de homomorfismos $\Psi^*(\mathcal{G}_\lambda)\to\mathcal{F}$, y el
límite inductivo de este sistema no es otro que $u^\sharp$.
\end{env}

\begin{env}[4.4.6]
\label{0.4.4.6-fr}
Sean $\mathcal{M}$, $\mathcal{N}$ dos $\mathcal{B}$-Módulos, $V$ un abierto
de $Y$, $U=\psi^{-1}(V)$; la aplicación $v\to\Psi^*(v)$ es un
homomorfismo
\[
  \operatorname{Hom}_{\mathcal{B}|V}(\mathcal{M}|V,\mathcal{N}|V)
  \longrightarrow
  \operatorname{Hom}_{\mathcal{A}|U}
  (\Psi^*(\mathcal{M})|U,\Psi^*(\mathcal{N})|U)
\]
\begin{sloppypar}
para las estructuras de $\Gamma(V,\mathcal{B})$-módulo
($\operatorname{Hom}_{\mathcal{A}|U}
(\Psi^*(\mathcal{M})|U,\allowbreak\Psi^*(\mathcal{N})|U)$ está normalmente provisto de una
estructura de $\Gamma(U,\psi^*(\mathcal{B}))$-módulo, y gracias al
homomorfismo
\oldpage[0\textsubscript{I}]{44}
canónico (3.7.2)
$\Gamma(V,\mathcal{B})\to\Gamma(U,\psi^*(\mathcal{B}))$, es, por tanto, también un
$\Gamma(V,\mathcal{B})$-módulo).
\end{sloppypar}
Se ve inmediatamente que estos homomorfismos son
compatibles con las restricciones, y por consiguiente definen un homomorfismo
canónico funtorial
\[
  \gamma:\mathcal{H}\!om_{\mathcal{B}}(\mathcal{M},\mathcal{N})
  \longrightarrow
  \Psi_*\bigl(\mathcal{H}\!om_{\mathcal{A}}
  (\Psi^*(\mathcal{M}),\Psi^*(\mathcal{N}))\bigr)~;
\]
a este homomorfismo le corresponde, por tanto, también el homomorfismo
\[
  \gamma^\sharp:
  \Psi^*\bigl(\mathcal{H}\!om_{\mathcal{B}}(\mathcal{M},\mathcal{N})\bigr)
  \longrightarrow
  \mathcal{H}\!om_{\mathcal{A}}
  (\Psi^*(\mathcal{M}),\Psi^*(\mathcal{N}))
\]
y estos homomorfismos canónicos son funtoriales en $\mathcal{M}$ y
$\mathcal{N}$.
\end{env}

\begin{env}[4.4.7]
\label{0.4.4.7-fr}
Supongamos que $\mathcal{F}$ (resp. $\mathcal{G}$) es una
$\mathcal{A}$-Álgebra (resp. una $\mathcal{B}$-Álgebra). Si
$u:\mathcal{G}\to\Psi_*(\mathcal{F})$ es un homomorfismo de
$\mathcal{B}$-Álgebras, $u^\sharp$ es un homomorfismo
$\Psi^*(\mathcal{G})\to\mathcal{F}$ de $\mathcal{A}$-Álgebras; esto resulta
de la conmutatividad del diagrama
\[
\begin{tikzcd}[column sep=huge,row sep=large]
\mathcal{G}\otimes_{\mathcal{B}}\mathcal{G}
  \arrow[r] \arrow[d]
  & \mathcal{G} \arrow[d,"u"] \\
\Psi_*(\mathcal{F}\otimes_{\mathcal{A}}\mathcal{F})
  \arrow[r]
  & \Psi_*(\mathcal{F})
\end{tikzcd}
\]
y de (4.4.4). Del mismo modo, si $v:\Psi^*(\mathcal{G})\to\mathcal{F}$ es un
homomorfismo de $\mathcal{A}$-Álgebras,
$v^\flat:\mathcal{G}\to\Psi_*(\mathcal{F})$ es un homomorfismo de
$\mathcal{B}$-Álgebras.
\end{env}

\begin{env}[4.4.8]
\label{0.4.4.8-fr}
Sean $(Z,\mathcal{C})$ un tercer espacio anillado,
$\Psi'=(\psi',\theta')$ un morfismo $(Y,\mathcal{B})\to(Z,\mathcal{C})$, y
$\Psi'':(X,\mathcal{A})\to(Z,\mathcal{C})$ el morfismo compuesto
$\Psi'\circ\Psi$. Sea $\mathcal{H}$ un $\mathcal{C}$-Módulo, $u'$ un
homomorfismo de $\mathcal{H}$ en $\mathcal{G}$; el compuesto
$v''=v\circ v'$ es por definición el homomorfismo de $\mathcal{H}$ en
$\mathcal{F}$ definido por
\[
  \mathcal{H}\xrightarrow{v'}\Psi'_*(\mathcal{G})
  \xrightarrow{\Psi'_*(v)}\Psi'_*(\Psi_*(\mathcal{F}))~;
\]
se verifica que ${v''}^\sharp$ es el homomorfismo
\[
  \Psi^*({\Psi'}^*(\mathcal{H}))
  \xrightarrow{\Psi^*({v'}^\sharp)}
  \Psi^*(\mathcal{G})
  \xrightarrow{v^\sharp}\mathcal{F}.
\]
\end{env}

